\documentclass[11pt]{article}

\usepackage[a4paper,margin=1in]{geometry}
\usepackage{amsmath,amssymb,amsfonts}
\usepackage{booktabs}
\usepackage{enumitem}
\usepackage{hyperref}
\usepackage{microtype}

\title{Contextual Chain: General Design Specification and Reference Profile\\
\large Version 2.0}
\author{Song-Ju Kim}
\date{August 22, 2026}

\begin{document}
\maketitle

\begin{abstract}
This document specifies \emph{Contextual Chain} as a general design framework in which
a participant maintains an evolving context and uses that context to guide subsequent
operational decisions. The term \emph{chain} refers to continuity between successive context
states; it does not require a blockchain, a particular ledger format, a fixed fork-choice rule, a
specific synchronization controller, or a cryptographic construction. Version 2.0 separates
the stable conceptual core from implementation-specific mechanisms. A recovery-oriented
event-driven ledger simulator is retained as a non-normative reference profile because it
provides the first concrete systems realization and reproducible evaluation environment.
Checkpointing, scoring, inconsistency estimation, synchronization adaptation, pruning, peer
selection, and cryptographic context binding are treated as replaceable modules rather than
defining features of Contextual Chain. The broader resource-asymmetry hypothesis---that a
participant following a shared context may track a narrow valid trajectory while a context-
uncertain observer may need to consider multiple plausible trajectories---remains a motivating
research direction, not a theorem or security claim of this specification.
\end{abstract}

\tableofcontents

\section{Status and Scope}
This document is Version 2.0 of the Contextual Chain design specification. The Version 1.0
record dated April 6, 2026 remains a historical description of the first operational realization.
Version 2.0 generalizes that formulation so that the design is not tied to one simulator, one
head-selection formula, one network model, or one resource-allocation policy.

Contextual Chain is a design family. An implementation conforms to the conceptual core
when it explicitly maintains an evolving operational context and uses that context in at least
one subsequent decision that affects acceptance, continuation, recovery, resource use, or security
state.

This document has four goals:
\begin{enumerate}[label=(\arabic*)]
  \item define the stable architectural ideas of Contextual Chain;
  \item identify interfaces that future realizations may instantiate in different ways;
  \item preserve the first recovery-oriented realization as a reproducible reference profile;
  \item state evidence and security boundaries without turning current experimental details into permanent protocol requirements.
\end{enumerate}

This document is not a standards-track interoperability specification and is not a complete
cryptographic security standard. It does not, by itself, claim Byzantine consensus, unforgeability,
Sybil resistance, post-quantum security, a formal memory lower bound, or bounded-memory
operation.

\section{Normative Language and Layering}
The words \textbf{MUST}, \textbf{SHOULD}, \textbf{MAY}, and \textbf{MUST NOT} apply only to the
conceptual interfaces and to an implementation that explicitly claims conformance to a named
reference profile. They do not require future Contextual Chain realizations to reuse the algorithms
or parameter values of the current simulator.

Version 2.0 distinguishes three layers:

\paragraph{Conceptual Core.}
Stable requirements that define the Contextual Chain design family.

\paragraph{Reference Profiles.}
Concrete, reproducible realizations of the core for a particular application domain. A profile
\textbf{MAY} define additional mandatory behavior for that profile only.

\paragraph{Experimental Extensions.}
Non-normative mechanisms used to investigate future design directions. They do not become
part of the conceptual core merely because they are evaluated.

\section{Core Conceptual Model}
\subsection{Context Sequence}
Let $C_t$ denote the context available to a participant after processing event or interaction $t$.
Contextual Chain treats operation as a sequence
\[
C_0 \rightarrow C_1 \rightarrow C_2 \rightarrow \cdots
\]
in which later decisions may depend on information carried forward from earlier accepted states.
The representation of $C_t$ is intentionally unspecified. It \textbf{MAY} contain a compact digest,
local graph state, recent observations, counters, scores, peer history, application state,
cryptographic commitments, or other information relevant to continuation. Different realizations
\textbf{MAY} retain different amounts of history.

\subsection{Context Update Interface}
A realization \textbf{MUST} define how context evolves. Abstractly,
\[
C_{t+1}=U(C_t,e_t),
\]
where $e_t$ is a newly observed event and $U$ is the realization-specific update rule. The update
\textbf{MAY} be deterministic, probabilistic, distributed, or cryptographically authenticated,
provided that the realization documents the semantics needed to reproduce or reason about its
behavior.

\subsection{Context-Informed Decision Interface}
A realization \textbf{MUST} use context in at least one consequential decision. Abstractly,
\[
a_t=D(C_t,x_t),
\]
where $x_t$ is a candidate input, observation, message, branch, peer, or action opportunity and
$a_t$ is an operational decision.

Examples include, but are not limited to:
\begin{itemize}
  \item accepting, rejecting, ranking, or deferring a candidate continuation;
  \item selecting a synchronization or recovery action;
  \item allocating communication, computation, storage, or verification effort;
  \item choosing peers or routes;
  \item deciding which state to retain, summarize, or prune;
  \item deriving or validating a context-bound authenticator.
\end{itemize}
No one item in this list is mandatory by itself.

\subsection{Continuity and Contextual Narrowing}
The general design intent is that a participant which has continuously followed a valid context
can use that continuity to narrow its next decision space. A participant lacking sufficient context
may face a wider set of plausible continuations. This principle is called \emph{contextual narrowing}
in this specification.

Contextual narrowing is a design objective, not a universal theorem. A particular realization
\textbf{MUST} state what information narrows the decision space and under what assumptions that
information remains useful.

\section{Core Design Requirements}
A realization claiming conformance to the Contextual Chain conceptual core \textbf{MUST} satisfy
the following requirements.

\paragraph{R1: Explicit evolving context.}
The realization \textbf{MUST} define an operational context $C_t$ whose content can change over time.

\paragraph{R2: Context dependence.}
At least one later operational decision \textbf{MUST} depend on $C_t$ or on a value derived from it.
A system that merely stores history but never uses it to influence later behavior is not a
Contextual Chain realization.

\paragraph{R3: Defined update semantics.}
The realization \textbf{MUST} define the events that update context and the conditions under which
an update is accepted, deferred, or discarded.

\paragraph{R4: Uncertainty behavior.}
The realization \textbf{MUST} define what happens when context is incomplete, inconsistent, stale,
or ambiguous. This \textbf{MAY} involve conservative behavior, multiple hypotheses, recovery,
explicit failure, or another documented policy.

\paragraph{R5: Recovery or continuation interface.}
The realization \textbf{MUST} define how operation resumes after a context gap, disagreement, or
reconnection when such events are within its application model. The mechanism is profile specific
and need not be gossip based.

\paragraph{R6: Modular context use.}
Context-dependent acceptance, synchronization, resource allocation, storage, and security
mechanisms \textbf{SHOULD} be separable where practical so that their effects can be evaluated
and replaced independently.

\paragraph{R7: Evidence boundary.}
A realization \textbf{MUST} distinguish architectural requirements from empirical findings and
from security assumptions. Experimental success of one module \textbf{MUST NOT} be presented
as proof that the module is universally required by Contextual Chain.

\section{Replaceable Module Interfaces}
The following modules are common but optional. A profile \textbf{MAY} instantiate any subset,
combine modules, or add new ones.

\subsection{Context Representation}
A profile \textbf{MAY} represent context as explicit state, summaries, hashes, commitments, recent
windows, learned embeddings, graph structure, or another representation. The conceptual core
does not prescribe a fixed encoding.

\subsection{Candidate Selection or Validation}
A profile \textbf{MAY} use context to rank or validate branches, messages, transactions, state
transitions, observations, or other candidates. The selection rule \textbf{MAY} be lexicographic,
score-based, probabilistic, rule-based, or cryptographic.

\subsection{Inconsistency or Confidence Estimation}
A profile \textbf{MAY} derive a local or distributed signal indicating uncertainty, inconsistency,
conflict, or confidence. The signal \textbf{MAY} use recent events, peer observations, statistical
estimators, model-based detectors, or other mechanisms.

\subsection{Synchronization and Recovery}
A profile \textbf{MAY} use periodic gossip, request/response repair, erasure coding, peer-assisted
reconciliation, state transfer, checkpoint retrieval, or another recovery mechanism. Context
\textbf{MAY} affect the timing, intensity, target, or content of synchronization, but adaptive
synchronization is not a mandatory feature of Contextual Chain.

\subsection{Resource Allocation}
A profile \textbf{MAY} use context to control communication, computation, storage, verification,
memory, or energy expenditure. The relevant resource metric \textbf{MUST} be defined by that
profile; simulator counters \textbf{MUST NOT} automatically be interpreted as physical bandwidth
or energy.

\subsection{Context Commitment and Authentication}
A future profile \textbf{MAY} bind context to cryptographic commitments, signatures,
challenge-response mechanisms, or other authenticators. Such a profile \textbf{MUST} define its
own security model and proof obligations. Operational context dependence alone is not
cryptographic authentication.

\subsection{Retention and Pruning}
A profile \textbf{MAY} retain full state, bounded windows, checkpoints, compressed summaries, or
pruned state. Bounded memory is therefore a profile property, not a requirement of the
conceptual core.

\section{Reference Profile R1: Intermittent-Network Recovery}
Reference Profile R1 preserves the first event-driven systems realization for reproducibility.
R1 is non-normative for the broader Contextual Chain design family; its rules are mandatory only
for an implementation claiming compatibility with R1.

R1 models a small ledger-like distributed system subject to delay, loss, temporary partition,
and rejoin. Each node maintains a local block graph, a current head, recent proposer and
inconsistency state, and a synchronization state. The simulator processes proposal, delivery,
and gossip-tick events.

The current R1 implementation contains the following mechanisms:
\begin{itemize}
  \item deterministic local head selection over visible tips;
  \item optional checkpoint metadata and proposer-based branch scoring;
  \item recent-window inconsistency estimation with hysteretic local state;
  \item periodic suffix-repair synchronization;
  \item fixed or context-triggered synchronization schedules;
  \item instrumentation for recovery, agreement, and synchronization-effort proxies.
\end{itemize}

The exact executable semantics, parameter defaults, evaluation variants, and accounting counters
are defined by the public reference implementation associated with this profile. A future R2,
R3, or domain-specific profile \textbf{MAY} use entirely different mechanisms while remaining
within the Contextual Chain conceptual core.

\subsection{Current R1 Interpretation}
The principal height-epoch configuration used in the current R1 evaluation treats checkpoint
level as an epoch label derived from height; in that configuration it should not be interpreted as
an independently validated contextual signal. Likewise, the current fuller head-selection formula
is a reference mechanism rather than a defining requirement of Contextual Chain.

R1 also includes a context-derived inconsistency signal that can influence synchronization
scheduling. The current native controller reads a simulator-wide aggregate and is therefore an
experimental oracle-assisted controller rather than a deployable decentralized estimator. Future
profiles \textbf{MAY} replace it with local, distributed, predictive, or application-specific
controllers.

\section{Evaluation Profiles and Mechanism Separation}
Contextual Chain does not prescribe one universal benchmark. A scientific evaluation
\textbf{SHOULD} separate mechanisms whenever feasible so that improvements attributed to
``context'' are not confounded with changes in resource quantity, scheduling, network conditions,
or unrelated control logic.

For R1, the public evaluation includes ablations of head-selection behavior and synchronization
policy, matched-seed schedule replay, and an equal-total-assigned-budget timing control. These are
evaluation instruments, not protocol primitives.

A future profile \textbf{SHOULD} document:
\begin{itemize}
  \item the context representation under test;
  \item which decisions depend on context;
  \item the relevant baseline that removes or neutralizes that dependence;
  \item the resource quantity held constant, if any;
  \item the source of randomness and matching procedure;
  \item the limits of any proxy metric.
\end{itemize}

\section{Current Empirical Status of R1 (Informative)}
Controlled experiments on the current R1 simulator support two qualitative conclusions. First,
synchronization provisioning is a major determinant of post-rejoin recovery in the studied
network model. Second, when total assigned synchronization opportunities are matched, a
context-triggered temporal allocation can outperform an approximately uniform temporal
allocation of the same assigned budget in the tested cases. Thus context can provide useful
control information about when recovery effort is applied.

The same controlled experiments do not establish a consistent independent advantage for the
current fuller head-selection formula over a height-only baseline. This negative result is a
property of the present R1 mechanism and configuration; it does not restrict future Contextual
Chain profiles from using other context-aware selection, validation, or commitment mechanisms.

Detailed numerical results, confidence intervals, seeds, schedules, code versions, and analysis
scripts belong to the corresponding reproducibility release and research article rather than to
the conceptual definition of Contextual Chain.

\section{Resource-Asymmetry Research Direction}
A broader motivation for Contextual Chain is the possibility of resource asymmetry created by
contextual continuity. Let an honest participant track one accepted trajectory $C_t$, while a
context-uncertain observer considers a set
\[
\mathcal{C}_t=\{C_t^{(1)},\ldots,C_t^{(m_t)}\}.
\]
If later verification or continuation depends on information carried by the accepted context,
then uncertainty may require maintaining, deriving, testing, or querying more alternatives.

This possibility can motivate designs in which the cost of uncertainty grows with the number,
depth, or diversity of plausible contexts. However, Version 2.0 deliberately leaves the cost
function, attacker model, challenge construction, and cryptographic primitive unspecified.
Different future realizations may pursue memory cost, communication cost, computation cost,
verification effort, search width, or another resource.

No lower bound or security theorem follows from the conceptual model alone.

\section{Security and Threat-Model Requirements}
A profile that makes security claims \textbf{MUST} state its threat model explicitly. Depending on
the application, this may include malicious proposers, Byzantine peers, equivocation, replay,
context withholding, state desynchronization, Sybil behavior, key compromise, or quantum-capable
adversaries.

The conceptual core does not assume that context is secret. A profile \textbf{MAY} use public
context, shared context, authenticated context, private context, or combinations thereof.

A cryptographic Contextual Chain profile \textbf{MUST} separately specify:
\begin{itemize}
  \item the security game or formal property;
  \item cryptographic assumptions and primitives;
  \item how context is bound to messages or state;
  \item adversarial capabilities and success criteria;
  \item soundness, completeness, or other required proofs.
\end{itemize}
The current R1 simulator does not satisfy these obligations and therefore makes no cryptographic
security claim.

\section{Implementation and Measurement Guidance}
Implementers \textbf{SHOULD} expose enough instrumentation to distinguish the following
categories when relevant:
\begin{itemize}
  \item policy-assigned opportunities or budgets;
  \item actions that survive filtering, loss, or eligibility checks;
  \item actual payload units scheduled or transmitted;
  \item physical bandwidth, latency, energy, or storage when measured in a real system.
\end{itemize}
These quantities \textbf{SHOULD NOT} be conflated.

A profile \textbf{SHOULD} also document whether its local state is retained indefinitely,
summarized, or pruned. If bounded memory or compact-state operation is claimed, the bound and
pruning semantics \textbf{MUST} be explicit.

\section{Versioning and Extensibility}
Contextual Chain versions are expected to evolve by adding or replacing reference profiles and
optional modules while preserving the conceptual core where possible.

Future versions \textbf{MAY}, among other directions:
\begin{itemize}
  \item introduce decentralized or predictive context estimators;
  \item replace gossip with other synchronization or reconciliation mechanisms;
  \item use topology-aware or learned peer selection;
  \item implement explicit pruning or compact-state retention;
  \item define context-bound authentication or authorization;
  \item instantiate the resource-asymmetry hypothesis with a formal adversary model;
  \item introduce cryptographic or post-quantum constructions;
  \item apply the framework outside ledger-like systems, including intermittent communications,
  sensing, distributed control, and other sequential decision systems.
\end{itemize}

A new mechanism \textbf{SHOULD} be described as part of the general Contextual Chain core only
if it is intended to be conceptually necessary across realizations. Otherwise it \textbf{SHOULD}
be introduced as an optional module, a reference-profile feature, or an experimental extension.

\section{Changes from Version 1.0}
Version 2.0 makes the following conceptual changes without invalidating Version 1.0 as a
historical implementation record:
\begin{enumerate}
  \item generalizes Contextual Chain from one operational ledger protocol to a design family based on evolving context and context-informed continuation;
  \item defines a small stable conceptual core and separates it from replaceable modules;
  \item moves checkpointing, fork choice, scoring, quarantine, and adaptive synchronization out of the universal definition and into the R1 reference profile;
  \item treats current controlled experiments as evidence about R1 rather than as permanent protocol requirements;
  \item preserves the resource-asymmetry idea as an open research direction without fixing one cryptographic or bookkeeping construction;
  \item makes explicit that future profiles may differ substantially in context representation, synchronization, storage, security, and application domain.
\end{enumerate}

\appendix
\section{Reference Profile R1 Reproducibility Note}
The exact R1 implementation is maintained as executable code rather than duplicated as a large
set of permanent normative equations in this general specification. The reproducibility release
records the simulator version, experiment runners, seeds, schedule controls, statistical analysis,
parameter provenance, and file checksums used for the associated controlled study.

For orientation only, the principal controlled R1 study uses a 20-node event-driven simulation,
a temporary partition followed by rejoin, noisy message delivery, and matched-seed comparisons
of alternative local-selection and synchronization policies. These choices define that study; they
do not constrain future Contextual Chain profiles.

\section{Profile Authoring Checklist}
A new Contextual Chain reference profile \textbf{SHOULD} document:
\begin{enumerate}
  \item the application and threat model;
  \item the representation of context $C_t$;
  \item the update rule $U$ and relevant event types;
  \item the context-informed decision function or policy $D$;
  \item uncertainty, inconsistency, and recovery behavior;
  \item optional synchronization, resource-allocation, security, and pruning modules;
  \item measurable outcomes and resource definitions;
  \item reproducibility artifacts and version identifiers;
  \item known limitations and claims that are explicitly out of scope.
\end{enumerate}

\section*{Acknowledgment of Scope}
This document is a general technical design specification and public disclosure. Its purpose is
to preserve the defining idea of Contextual Chain while allowing future realizations to evolve.
The current recovery simulator is one reference profile, not the endpoint of the design. The
specification therefore separates stable conceptual requirements from implementation-specific
mechanisms, empirical results, and future security constructions.

\end{document}
